\chapter{Cerebral Visual Impairment}
\index{Cerebral Visual Impairment}
\label{sec:Cerebral Visual Impairment}

\section{Overview}
\begin{itemize}[noitemsep]
  \item Cerebral (or cortical) visual impairment (CVI) is reduced vision caused by damage to the brain's visual processing areas rather than to the eyes themselves
  \item It is the leading cause of visual impairment in children in developed countries
  \item Vision loss ranges from mild to severe and often coexists with other neurological conditions
\end{itemize}
\section{Causes}
This condition happens because of damage to the visual pathways or visual cortex of the brain. In children, common causes include prematurity and perinatal hypoxia, periventricular leukomalacia, intraventricular hemorrhage, infection during pregnancy (e.g., cytomegalovirus, toxoplasmosis), hydrocephalus, and traumatic brain injury. In adults, stroke, head trauma, infection, and brain tumors are the main causes.
\section{Risk Factors}
\begin{itemize}[noitemsep]
  \item Premature birth and low birth weight
  \item Hypoxic-ischemic encephalopathy at birth
  \item Hydrocephalus and shunted cerebrospinal fluid
  \item Seizures (epilepsy)
  \item Cerebral palsy and developmental delay
  \item Stroke, trauma, or brain tumors in adults
\end{itemize}
\section{Signs and Symptoms}
\subsection{Common symptoms}
\begin{itemize}[noitemsep]
  \item Visual acuity poorer than expected for the eye exam findings
  \item Difficulty recognizing faces or objects
  \item Preferential attention to moving or brightly colored objects
  \item Difficulty seeing in cluttered or visually busy environments
  \item Variable vision that changes with fatigue, light, or novelty
  \item Reduced visual fields, often lower visual field loss
\end{itemize}
\subsection{Emergency warning signs}
\begin{itemize}[noitemsep]
  \item Sudden loss of vision with other neurological symptoms such as weakness, confusion, or headache requires urgent evaluation
\end{itemize}
\section{Progression and Stages}
CVI is not progressive when caused by a static brain injury; however, the functional vision often improves over time, especially in children, as the brain matures and with rehabilitation. In adults with stroke or degenerative disease, the course depends on the underlying cause. Recovery can continue for months to years after injury.
\section{Complications}
\begin{itemize}[noitemsep]
  \item Delayed motor, cognitive, and language development in children
  \item Difficulty with learning and school performance
  \item Increased risk of falls and accidents
  \item Social isolation and behavioral challenges
  \item Overlapping sensory and neurological disabilities
\end{itemize}
\section{Diagnosis}
\begin{itemize}[noitemsep]
  \item Ophthalmological examination showing normal or near-normal eye structure
  \item Neuroimaging (MRI) to identify brain lesions
  \item Visual evoked potentials (VEP) to assess visual pathway function
  \item Functional vision assessment by a specialist team
  \item EEG and neurological evaluation for associated seizures
\end{itemize}
\section{Prevention}
\begin{itemize}[noitemsep]
  \item Preventing prematurity and optimizing neonatal care
  \item Prompt treatment of hydrocephalus and seizures
  \item Prevention of head injury and infection
  \item Stroke prevention measures in adults (blood pressure, cholesterol, healthy lifestyle)
\end{itemize}
\section{Treatment Options}
\begin{itemize}[noitemsep]
  \item Early intervention and developmental programs for children
  \item Specialized educational support and visual stimulation training
  \item Environmental modifications: reducing clutter, enhancing contrast and lighting
  \item Multidisciplinary rehabilitation (vision, occupational, and physical therapy)
  \item Treatment of the underlying neurological condition
  \item Assistive technology and accommodations in school and at home
\end{itemize}
\section{Living with It}
\begin{itemize}[noitemsep]
  \item Provide predictable, simplified visual environments
  \item Use familiar objects, routines, and verbal cues
  \item Involve teachers, therapists, and family in a consistent approach
  \item Regularly reassess functional vision, as it can improve over time
  \item Seek support networks for families affected by CVI
\end{itemize}
\section{Outlook}
The outlook is variable. Many children with CVI make significant visual gains with early intervention and support, particularly when the underlying injury is static. Functional vision and quality of life often improve substantially with rehabilitation and appropriate educational strategies.
